% ===========================================================================
%  正文迁移块（由 main.tex 在 appendices 内 \input 引入）
%
%  竞赛格式规范要求正文不超过 30 页，而问题四新增了整整一节。按"无损优先"
%  的原则，把问题一至问题三中**过程性、细化**的分析整块移到此处，正文原位
%  只保留结论段并指向本附录。搬运方式是"剪切 + 保留标签"，所有图号表号由
%  LaTeX 重新编号，正文的 \ref 自动解析到附录，内容一字未删。
% ===========================================================================

\section{正文迁移的过程性分析与对照明细}
\label{app:moved}

本附录收录从正文移出的过程性内容。每节开头注明它在正文中的原位，正文对应
位置留有含关键数字的结论段；图号表号由 \LaTeX{} 统一编号，正文的引用会
自动指向本附录，\textbf{数据与推导过程一字未删}。

\subsection{问题二：紧急购电的触发机制（原正文 6.9.1 节）}

图~\ref{fig:p2-forecast} 给出四个指定日期的净负荷预测、风险余量与实际净负荷。
三个层次的关系清晰可见：净负荷预测（虚线）给出基准，风险余量
$Q_{\alpha}$（橙色阴影）在预测之上叠加一段风险修正量，得到用于日前求解的
风险净负荷（橙实线），而实际净负荷（深蓝实线）则在其上下波动。

需要说明，$Q_{\alpha}$ 是历史误差的经验分位数，\textbf{并不保证为正}：正式
区间 $334\times144=48096$ 个时段中有 $4069$ 个（$8.5\%$）为负，最小值
$-40.6$~kWh、最大值 $258.5$~kWh。这些负值出现在误差分布整体偏负（长期被
高估）的时段，此时"分位数修正"正确地指示了\textbf{可以少买}——它与总体上
"风险余量抬高计划量"的方向并不矛盾，因为决定费用的始终是分位数本身，而不是
"加一个正数"这一形式。

关键在于\textbf{实际曲线落在风险净负荷之上的部分}。图~\ref{fig:p2-forecast}
用倒三角标出了发生紧急购电的时段，$9$ 月 $23$ 日出现了 $19$ 段紧急购电，
$3$ 月 $20$ 日 $1$ 段，其余两个日期则全天为零。但逐时段核对可以发现，
\textbf{超出风险净负荷只是发生紧急购电的必要条件，而不是充分条件}：

\begin{itemize}
  \item $6$ 月 $21$ 日几乎被风险净负荷\textbf{完整包住}（$144$ 段中仅 $1$ 段超出，
        超出量 $4.7$~kWh）；
  \item $3$ 月 $20$ 日与 $12$ 月 $21$ 日分别有 $50$ 段与 $18$ 段超出风险净负荷，
        但最大超出仅 $79.9$~kWh 与 $84.4$~kWh，且都出现在储能尚有余量的时段，
        被电池吸收，故\textbf{全天零紧急购电}；
  \item $9$ 月 $23$ 日有 $71$ 段超出，最大超出 $133.7$~kWh，且\textbf{当日储能
        在 $20{:}40$ 已被放到下限 $1200$~kWh，并在其后 $9$ 个时段钉在下限上}；
        电池一旦触底，其后每一段缺口都只能按 $5$ 倍电价补购。
\end{itemize}

这解释了 $9$ 月 $23$ 日晚间的 $19$ 段紧急购电：其中 $15$ 段的实际净负荷高于
风险净负荷，另 $4$ 段（$19{:}00$、$19{:}30$、$19{:}40$、$20{:}30$）的实际值
甚至\textbf{低于}风险净负荷 $9.8$--$30.5$~kWh，却仍分别紧急购电
$51.0$--$71.0$~kWh——因为电池已无电可放。可见真正的触发
机制是\textbf{"实际净负荷超出风险余量的部分"与"储能剩余可放能力"两者之差}：
前者为正只是故事的开始，只有当储能无法再兜底时，才转化为 $5$ 倍电价的支出。
这三个日期的差别正在于此——$3$ 月 $20$ 日与 $12$ 月 $21$ 日的超出量小
（超幅超过 $50$~kWh 的时段各仅 $4$ 个与 $2$ 个），$9$ 月 $23$ 日则达 $22$ 个。

由此得到本文最重要的定性结论：\textbf{紧急购电成本来自"预测偏低"，而非
"预测不准"}——预测偏高只会造成弃电（无直接费用），预测偏低才会在储能耗尽时
触发 $5$ 倍电价。这也是选择\textbf{分位数}而非均值作为计划依据的根本原因。

\begin{figure}[htbp]
  \centering
  \includegraphics[width=.66\textwidth]{p2_fig1_forecast.png}
  \caption{指定日期的净负荷预测、风险余量与实际净负荷}
  \label{fig:p2-forecast}
\end{figure}

\subsection{问题二：预测误差对储能轨迹的扰动（原正文 6.9.3 节）}

图~\ref{fig:p2-storage} 对比了四个指定日期的日前参考储电量 $\bar E$ 与实际
执行后的储电量 $E$。两者的差异正是预测误差在执行层的体现：参考轨迹是 $0{:}00
$ 基于风险净负荷算出的最优路径，而实际轨迹则因实际负载与光伏的偏离而上下
浮动。图中同时标出了 $[1200,10800]$~kWh 的运行边界与 $E^{\mathrm{tar}}=6000$
~kWh 的参考水平。可以看到实际轨迹始终未越界，且四个日期的日初与日末储电量
\textbf{各不相同}（例如 $3$ 月 $20$ 日 $1982.53\to1624.64$~kWh，$6$ 月 $21$
日 $4554.68\to4903.38$~kWh），这正是问题二\textbf{跨日连续、不要求日循环}
的直接证据——若强行要求每天回到 $6000$~kWh，反而会人为限制模型的调节能力。

\begin{figure}[htbp]
  \centering
  \includegraphics[width=.66\textwidth]{p2_fig3_storage.png}
  \caption{指定日期的日前参考储电量与实际储电量}
  \label{fig:p2-storage}
\end{figure}

\subsection{问题二：全年费用分解与对照策略（原正文 6.9.4 节）}

图~\ref{fig:p2-yearly} 从全年尺度展示结果。
图~\ref{fig:p2-yearly}(a) 给出本文策略的逐月费用分解。月费用最高的三个月是
$12$ 月（$175.9$ 万元）、$6$ 月（$155.2$ 万元）与 $7$ 月（$154.1$ 万元），
最低的三个月是 $4$ 月（$96.6$ 万元）、$5$ 月（$100.4$ 万元）与 $10$ 月
（$112.5$ 万元），最高月约为最低月的 $1.8$ 倍。这一形态主要由负载与光伏的
季节错配决定：$12$ 月光伏出力全年最弱而负载处于高位，$6$--$7$ 月则因制冷
负荷抬高用电水平，而春季与初秋光伏尚可、负载适中，故费用最低。

值得注意的是\textbf{月度总费用的排序与紧急购电费的排序并不一致}：紧急购电费
在 $6$ 月达到 $190493$ 元，为全年最高，\textbf{占全年紧急购电费的 $25.2\%$}，
$12$ 月（$123830$ 元）与 $9$ 月（$72567$ 元）次之；而 $11$ 月的紧急购电费
仅 $11265$ 元、为全年倒数第二低，其总费用却仍排在第 $5$ 位。这说明\textbf{月费用
的高低主要仍由负载水平决定，紧急购电费则取决于预测偏差恰好与储能容量紧张
相遇的频率}，二者是两个不同的驱动因素。

为量化风险余量的价值，构造三组对照策略，全部使用\textbf{同一预测器、同一
实时执行规则、同一评价区间（$2$--$12$ 月共 $334$ 天）}，并各自单独标定参数：

\begin{itemize}
  \item \textbf{预测均值策略}：令 $\alpha=\text{None}$，即完全不加风险余量，
        直接用点预测制定日前计划，其余规则完全相同；
  \item \textbf{无储能策略}：仍提前制定购电计划、仍按 $5$ 倍价格补缺，但
        不启用储能，用于单独评价电池的价值；
  \item \textbf{事后理想策略}：用当天\textbf{实际}净负荷替换预测值求解同一
        日前 MILP，代表"完全预知当天净负荷"的费用参照。须注意它\textbf{每天
        从主策略的实际日初储电量出发、却不承担日末留存义务}，等于每天白拿一次
        主策略已付费的库存电量（$334$ 天累计 $306254$~kWh，按平均电价折合约
        $23.5$~万元），同时它在制定当天计划时已看到当天实际净负荷。因此它是
        一个\textbf{可执行性不成立的乐观下界}，而非可达方案；即使按库存口径
        修正（库存单价取平均电价至其可实现边际价值上界 $\eta_c\eta_d\max_t p_t$
        之间），它与主策略的差距仍有 $10.6\%$--$11.4\%$。
\end{itemize}

表~\ref{tab:p2-strategy} 给出四者的费用对比。本文风险修正策略全年合计
$13939720$~元，较预测均值策略的 $14903679$~元节省 $963959$~元、较无储能
策略的 $18519391$~元节省 $4579670$~元，分别相当于\textbf{各基准费用的}
$6.47\%$ 与 $24.73\%$（表~\ref{tab:p2-strategy} 末列给出同一差额，但以本模型
费用为分母，见该表注）。事后理想
策略为 $12110855$~元，比本文策略低 $1828866$~元（$13.1\%$）。但该值
\textbf{偏乐观}：理想策略每天直接动用主策略已按计划电价付过费的库存电量，
且不要求日末留存——其 $334$ 天的日末储电量\textbf{全部停在 $1200$~kWh 的下限}，
正是这一"白拿库存"的直接证据。因此 $13.1\%$ 是高估的上限，而非预测不确定性
带来的不可避免代价；本文策略由此仍有明确的改进空间。

更值得注意的是费用结构的变化：预测均值策略的\textbf{紧急购电费高达
$2708867$~元，是本文策略 $755499$~元的 $3.6$ 倍}，而其计划购电费反而更低
（$12194812$~元 对 $13184221$~元）。这正是"代价不对称"的直接体现——不加
风险余量时，日前计划买的电更少、看起来更省，但代价是运行时频繁触发 $5$ 倍
电价的紧急购电，总费用反而更高。\textbf{风险余量本质上是用"少量、分散、
可预期的计划购电溢价"替换"大量、集中、代价高昂的紧急购电"}。

同时应注意，四个策略的\textbf{期末储电量并不相同}（本文策略
$1559.7$~kWh，预测均值策略 $5960.7$~kWh，无储能策略 $6000.0$~kWh）。
期末储电量差异意味着库存并未被完全对等地消耗，因此本文同时报告该指标，
避免把"多消耗库存"误当作算法收益。

% 两张全年尺度的图合并为同一浮动体（共享标题、节约一页篇幅）
\begin{figure}[htbp]
  \centering
  \begin{subfigure}{.76\textwidth}
    \centering
    \includegraphics[width=\textwidth]{p2_fig4_monthly.png}
    \caption{逐月费用分解与策略费用对比}
  \end{subfigure}\\[3pt]
  \begin{subfigure}{.76\textwidth}
    \centering
    \includegraphics[width=\textwidth]{p2_fig5_cumulative.png}
    \caption{各策略累计费用与跨日储电量}
  \end{subfigure}
  \caption{全年费用分解与策略对照}
  \label{fig:p2-yearly}
\end{figure}

\subsection{问题二：弃用电量的成因（原正文 6.9.6 节）}

全年弃用电量 $2463544.98$~kWh，占同期光伏发电量 $19068890.0$~kWh 的 $12.9\%$。
题目未对弃电设置惩罚或收益，故弃电不带来额外扣款，但它仍是日前计划与实时
运行之间偏差的\textbf{指示器}。逐时段分解显示弃电来自\textbf{两个性质不同}的
来源：$71.8\%$（$1769824$~kWh）是光伏富余超出储能吸纳能力（储能受 $5000$~kW
功率上限与 $10800$~kWh 容量上限双重约束）；另 $28.2\%$（$693721$~kWh）发生在
\textbf{光伏并无富余}的时段，只能来自\textbf{已按计划电价付费、却因实际负载
低于预期而用不掉的计划购电量}，其当初购电成本 $466267$~元约占全年计划购电费
的 $3.54\%$——即约三成弃电是"按 $5$ 倍电价买保险"预付而未出险的保费。
$\alpha$ 敏感性验证表明这一
取舍\textbf{经过权衡且最优}：$\alpha$ 由 $0.80$ 降至 $0.50$ 时弃电由 $2628.6$
降至 $1493.7$~MWh，但紧急购电量由 $177.9$ 猛增至 $636.1$~MWh，
\textbf{总费用反由 $1403.8$ 万元升至 $1508.8$ 万元}。此外，紧急购电负担在时间上
\textbf{高度集中}：$334$ 天中 $143$ 天发生紧急购电，而 $6$ 月 $1$--$2$ 日两天
即占全年紧急购电费的 $13.1\%$，说明它由少数连续多日的高压时段主导。

\begin{figure}[htbp]
  \centering
  \includegraphics[width=.66\textwidth]{p2_fig6_daily.png}
  \caption{紧急购电量的季节分布与逐日费用构成}
  \label{fig:p2-daily}
\end{figure}

\subsection{问题一：购电策略的逐段时序特征（原正文 5.5.1 节）}

图~\ref{fig:p1-purchase} 给出全天逐时段的计划购电量与电价的对照关系。可以
清晰地看到购电行为被\textbf{严格地压缩到了低电价时段}：

\begin{itemize}
  \item 在 $0{:}00$--$6{:}00$ 的夜间低谷（电价约 $0.37$--$0.91$~元/kWh），
        购电量维持在每 $10$ 分钟 $1400$~kWh 上下，其中 $0{:}00$--$0{:}40$
        因需同时供给负载并为储能充电而达到全天最大的
        $1406.6$--$1409.8$~kWh/10\,min；
  \item 在 $5{:}30$--$5{:}40$ 电价降至全天最低（$0.3713$~元/kWh）时，购电量
        为 $1383.66$~kWh/10\,min，仍处高位；此后随储能充满
        （$5{:}40$ 触及上限 $10800$~kWh）迅速退出购电，$5{:}50$ 起连续
        $31$ 个时段购电量降为零；
  \item 在 $7{:}50$--$9{:}10$ 与 $17{:}40$--$20{:}50$ 两个电价高峰段
        （电价最高达 $1.3952$~元/kWh，约为低谷的 $3.76$ 倍），购电行为被
        显著抑制：早高峰 $7{:}30$--$10{:}00$ 的购电量\textbf{恰好为零}，
        完全由储能放电支撑；晚高峰 $17{:}00$--$21{:}00$ 购电量回落至
        $0$--$826.08$~kWh/10\,min（$24$ 个时段中仅 $12$ 个时段有购电），
        降幅虽大但仍保留了部分购电——因为此时储能储电量已接近下限，
        继续放电将破坏后续时段的供电能力；
  \item 在光伏出力超过小区负载的 $9{:}30$--$14{:}20$（共 $30$ 个时段）内，
        购电量在 $22$ 个时段\textbf{降为零}，微网完全由光伏供电并有盈余为
        储能充电；其余 $8$ 个时段（$11{:}00$--$12{:}30$）仍有
        $179.28$--$575.17$~kWh/10\,min 的购电。这并非供电需要，而是模型
        捕捉到了正午时分电价出现的局部低谷（$0.4113$--$0.5007$~元/kWh，
        低于相邻时段）：此时储能尚有充足容量（储电量由 $2971$~kWh 升至
        $9144$~kWh），故在光伏已经富余的情况下仍提前购入低价电存入储能，
        留待傍晚 $1.3952$~元/kWh 的高价时段使用。这正是
        表~\ref{tab:p1-buy} 中 $10{:}00$、$14{:}00$、$20{:}00$ 三个时段购电量
        为零、而 $12{:}00$、$16{:}00$、$18{:}00$ 仍保留一部分购电的原因
        ——$16{:}00$、$18{:}00$ 的购电则是因储电量已达上限 $10800$~kWh、
        光伏出力已低于负载，只能由外网补足。
\end{itemize}

\begin{figure}[htbp]
  \centering
  \includegraphics[width=.66\textwidth]{p1_fig2_purchase.png}
  \caption{问题一：全天计划购电量与电价}
  \label{fig:p1-purchase}
\end{figure}

\subsection{问题一：储能运行的逐段细节（原正文 5.5.2 节）}

图~\ref{fig:p1-storage} 给出储能设备的充放电功率与储电量轨迹，呈现出典型的
两充两放格局：

\begin{enumerate}
  \item \textbf{第一充放循环}：$0{:}00$--$5{:}40$ 利用夜间低谷电价充电，
        储电量由 $6000$~kWh 升至 $10800$~kWh（$5{:}40$ 首次触及上界）；
        $5{:}50$--$9{:}30$ 在早高峰放电，储电量一路降至全天次低点
        $1834.62$~kWh；
  \item \textbf{第二充放循环}：$9{:}30$--$16{:}00$ 吸收光伏盈余并在正午电价
        局部低谷时补充购电充电（储电量自 $1834.62$~kWh 回升，
        $13{:}50$ 突破 $10500$~kWh，$16{:}00$ 再次达到上界
        $10800$~kWh）；$18{:}00$--$20{:}40$ 在晚高峰放电，
        储电量降至下界 $1200$~kWh（$20{:}40$ 首次触及，并保持至
        $21{:}50$）；
  \item $21{:}50$--$24{:}00$ 再次充电，将储电量精确送回 $6000$~kWh，满足
        日循环条件。
\end{enumerate}

\begin{figure}[htbp]
  \centering
  \includegraphics[width=.66\textwidth]{p1_fig3_storage.png}
  \caption{问题一：储能设备充放电功率与储电量}
  \label{fig:p1-storage}
\end{figure}

\subsection{问题一：全天电量平衡与弃光（原正文 5.5.3 节）}

图~\ref{fig:p1-balance} 给出全天电量的供给与需求构成。供给侧为光伏发电
$55482.84$~kWh、电网购电 $59482.70$~kWh 与储能放电 $16799.94$~kWh，合计
$131765.48$~kWh；需求侧为小区负载 $111024.81$~kWh 与储能充电 $20740.67$~kWh，
合计亦为 $131765.48$~kWh，两侧严格相等。

一个值得注意的结果是：本模型的\textbf{弃光电量为零}，即 $s_t^* = 0$ 对所有
$t$ 成立。储能设备将光伏出力超过负载的全部盈余电量（$6247.96$~kWh）尽数
吸纳，未造成任何浪费。这是 $10$ 分钟精细调度粒度带来的直接收益——若采用
粗粒度调度，光伏盈余与储能充电能力在时间上难以精确匹配，将不可避免地产生
弃光（详见正文第 \ref{sec:sensitivity} 节）。

\begin{figure}[htbp]
  \centering
  \includegraphics[width=.66\textwidth]{p1_fig4_balance.png}
  \caption{问题一：全天电量平衡构成}
  \label{fig:p1-balance}
\end{figure}

\subsection{问题三：预报使用的完整边际分析（原正文 7.7 节）}
\label{app:p3-combo-detail}

表~\ref{tab:p3-combo} 与图~\ref{fig:p3-combo} 给出八种预报使用方式的实际总费用。
为回答"是否需要其他时刻的预报"，还要看\textbf{边际}改善
$\Delta_h(S)=C(S)-C(S\cup\{h\})$，即已经用了 $S$ 之后再加用时刻 $h$ 能省多少。
其中 $S=\varnothing$ 一列表示"只用 $0{:}00$ 的预报、关闭全天调整"的对照。结果
如下（单位：万元）：

\begin{center}
\footnotesize
\begin{tabular}{lrrrr}
  \toprule
  已有组合 $S$ & $\Delta_{6}(S)$ & $\Delta_{12}(S)$ & $\Delta_{18}(S)$ & 参考 \\
  \midrule
  $\varnothing$（仅 $0{:}00$） & $49.78$ & $94.59$ & $\mathbf{117.82}$ & 单独引入 \\
  $\{6\}$        & ---     & $59.91$ & $\mathbf{94.69}$ & \\
  $\{12\}$       & $15.10$ & ---     & $\mathbf{45.17}$ & \\
  $\{18\}$       & $\mathbf{26.66}$ & $21.94$ & --- & \\
  $\{6,12\}$     & ---     & ---     & $46.23$ & \\
  $\{6,18\}$     & ---     & $11.44$ & --- & \\
  $\{12,18\}$    & $16.16$ & ---     & --- & \\
  \bottomrule
\end{tabular}
\end{center}

由此可以给出三个层次明确的结论。

\textbf{其一，$18{:}00$ 的预报最值得引入，且在任何已有组合下都不可替代。}单独
引入时它是三者中改善最大的（省 $117.82$ 万元）；在 $\{6\}$、$\{12\}$、
$\{6,12\}$ 之上再引入它仍分别省 $94.69$、$45.17$、$46.23$ 万元，都是在当时
所有可选时刻中最大的。值得强调的是，\textbf{这个排序与预报精度无关}：按
表~\ref{tab:p3-skill}，在四版共同覆盖的白天窗口 $12{:}00$--$18{:}00$ 上
$12{:}00$ 版最准（平均绝对误差 $181.3$~kW，优于 $6{:}00$ 版的 $304.9$~kW 与
$0{:}00$ 版的 $445.8$~kW），而 $18{:}00$ 版\textbf{对这一窗口没有任何覆盖}。
$18{:}00$ 版的价值来自它\textbf{决定的是哪一个交付块}：它一旦提交就锁定晚间
$18{:}00$--$24{:}00$，这一段既覆盖电价高峰、又是日末回充窗口，而且是全天唯一
\textbf{没有光伏}的时段——风险全部落在没有外部预报、只能靠同星期往期外推的
负载上。换言之，一份预报的价值取决于\textbf{它所覆盖的时段有多难预测、代价有
多大}，而不只取决于它自己有多准。

\textbf{其二，$12{:}00$ 的边际价值在已有相邻版本后迅速衰减。}它在
$\{18\}$ 之上只省 $21.94$ 万元、在 $\{6,18\}$ 之上只省 $11.44$ 万元，远小于
它单独引入时的 $94.59$ 万元。原因是它的信息与相邻两个时刻\textbf{高度重复}：
它对当天下午的修正，$6{:}00$ 已经预报过一次，$18{:}00$ 又会覆盖相同的目标时段
（三个版本的提交块首尾相接，$12{:}00$ 版负责的第 3 块被前后两个版本"夹"在中间）。

\textbf{其三，"哪个版本更值钱"依赖于已经拥有哪些版本。}单独引入时，$12{:}00$
优于 $6{:}00$（$94.59$ 对 $49.78$ 万元）；但在已有 $\{18\}$ 之后次序\textbf{反转}
——此时加 $6{:}00$ 省 $26.66$ 万元，加 $12{:}00$ 只省 $21.94$ 万元。这与
$\Delta_6(\{6,12\})=15.10$ 万元同源：$6{:}00$ 版在 $\{12\}$ 之上仍能省
$15.10$ 万元，说明它带来的不只是"第 2 块的预报"，还包括以当日实际储电量和上午
已实现的光伏为初值、重新求解整个剩余时段所换来的好处；而这一好处在 $12{:}00$
版上已经被 $18{:}00$ 版部分地实现了。因此\textbf{边际价值不能脱离已有信息集来
排序}。

\textbf{若必须精简。}按"撤下的代价"排序：撤掉 $18{:}00$ 版全年多花
$462{,}252.98$ 元、撤掉 $6{:}00$ 版多花 $161{,}579.57$ 元、撤掉 $12{:}00$ 版
只多花 $114{,}388.95$ 元。故三个时刻\textbf{全部引入}是最省的；若为了少一次
订阅与一次求解而必须舍去一个，舍去 $12{:}00$ 版的代价最小，约为全年费用的
$0.85\%$。

\textbf{一个必须说明的反常现象：仅用 $0{:}00$ 预报并不比问题二更省。}对照方式
（只用 $0{:}00$ 版预报、关闭全天调整）全年费用为 $14{,}977{,}809.03$ 元，比
问题二的历史外推策略 $13{,}939{,}720.46$ 元\textbf{高出 $1{,}038{,}088.57$ 元}。
需要注意这里\textbf{不能}推出"附件 3 的预报不如历史外推"：该对照同时切换了
两件事——把计划所用的光伏预测换成 $0{:}00$ 版预报，并且\textbf{关闭了全部日内
调整}。它的紧急费高达 $2{,}382{,}672.91$ 元（占自身费用的 $15.9\%$），而问题二
的紧急费只有 $755{,}499.17$ 元，差距主要来自"不做任何日内调整"这一项。这个数字
只说明\textbf{"仅凭一次 $0{:}00$ 预报、不做日内调整"不是一个够用的方案}，不足以
单独论证预报质量。

\textbf{信息收益与机制收益不做归因拆分。}上面格上的每一条覆盖边（例如
$\{6\}\to\{6,12\}$ 省 $59.91$ 万元、$\{6,12\}\to\{6,12,18\}$ 省 $46.23$ 万元）
都是在\textbf{机制不变}的前提下只更换预报集合，因此它们可以直接读作\textbf{纯
信息增量}。但从对照 $\varnothing$ 到全集的那 $155.92$ 万元的差额同时包含"是否
启用日内调整机制"与"新发布预报的信息"两部分，本文\textbf{不对它作进一步拆分}：
严格分离需要额外设计一个"在三个发布时刻重新求解、但继续沿用 $0{:}00$ 版预报"
的对照回放，这已超出本题的评价口径。因此正文中 $10.41\%$ 的改善应理解为
\textbf{滚动调整机制与新预报信息共同的贡献}，而各版本之间的横向比较才是纯粹的
信息价值。


\subsection{问题四：八种预报组合的完整边际分析（原正文 8.6 节）}
\label{app:p4-combo-detail}

记 $\Delta_h(S)=J(S)-J(S\cup\{h\})$ 为在已有预报集合 $S$ 之上\textbf{再订一版}
$h$ 所省下的费用。八种组合各自独立标定后完整回放 334 天，$12$ 条"多用一个
时刻"的覆盖边上的增量全部列于表~\ref{tab:p4-combo-detail}。

% 本文件由 03_代码/p4_tables.py 自动生成，请勿手工修改。
% 数据源：06_支撑材料/p4_results.json

\begin{table}[htbp]
  \centering
  \caption{4-3 八种预报组合的边际价值（波动电价，各组合独立标定；单位：元）}
  \label{tab:p4-combo-detail}
  \footnotesize
  \begin{tabular}{llrr}
    \toprule
    已有预报 $S$ & 再加一版 $h$ & 增量 $\Delta_h(S)$ & 折合 \\
    \midrule
    $\varnothing$ & $6{:}00$ & 356,384 & 35.64 万元 \\
    $\varnothing$ & $12{:}00$ & 286,251 & 28.63 万元 \\
    $\varnothing$ & $18{:}00$ & 522,757 & 52.28 万元 \\
    \midrule
    $\{6\}$ & $12{:}00$ & 147,030 & 14.70 万元 \\
    $\{6\}$ & $18{:}00$ & 326,395 & 32.64 万元 \\
    $\{12\}$ & $6{:}00$ & 217,163 & 21.72 万元 \\
    $\{12\}$ & $18{:}00$ & 193,552 & 19.36 万元 \\
    $\{18\}$ & $6{:}00$ & 160,022 & 16.00 万元 \\
    $\{18\}$ & $12{:}00$ & -42,954 & -4.30 万元 \\
    \midrule
    $\{6,12\}$ & $18{:}00$ & 117,083 & 11.71 万元 \\
    $\{6,18\}$ & $12{:}00$ & \textbf{-62,282} & \textbf{-6.23 万元} \\
    $\{12,18\}$ & $6{:}00$ & 140,694 & 14.07 万元 \\
    \bottomrule
  \end{tabular}
  \par\vspace{2pt}
  \begin{minipage}{.92\textwidth}\footnotesize
    注：$\Delta_h(S)=J(S)-J(S\cup\{h\})$ 为在已有集合 $S$ 之上再订一版 $h$ 所省下的费用。所有组合都在同一条附件 4 实际价格路径上结算，且各自重新标定 $\alpha,\rho,\lambda$，因此增量可以横向比较。加粗行是\textbf{边际价值最低}的一条；表中\textbf{增量为负}的行意味着在已有预报之上再订一版\textbf{反而更贵}——该现象及其成因见附录该节说明。
  \end{minipage}
\end{table}


两点值得强调。其一是\textbf{边际递减的方向并不对称}。$12{:}00$ 一版的价值在
不同起点下\textbf{连符号都会变}：什么都没有时它值 $28.63$ 万元，而已有
$18{:}00$ 时\textbf{转为 $-4.30$ 万元}，已有 $\{6{:}00,18{:}00\}$ 时是
$-6.23$ 万元。反过来，$6{:}00$ 与 $18{:}00$ 之间没有这种塌缩——$6{:}00$
在已有 $\{12{:}00,18{:}00\}$ 时仍值 $14.07$ 万元，$18{:}00$ 在已有
$\{6{:}00,12{:}00\}$ 时仍值 $11.71$ 万元。这说明\textbf{$12{:}00$ 与
$18{:}00$ 携带的信息高度重叠}（两者之间只隔 6 小时，而两者相隔的时段恰是光伏
出力已被 $12{:}00$ 版确认、晚间电价高峰尚未到来的区间），而 $6{:}00$ 版覆盖的是
\textbf{另一类信息}：它是三个时刻中最靠近凌晨与上午的，对上午的量价判断最有用。
因此"只多订一版"应选 $18{:}00$；若已订 $18{:}00$，加 $12{:}00$ 已无信息增量，
不如把这一版留给 $6{:}00$。

那两条负增量\textbf{不应读作"多一版预报有害"}。它们恰是"在已有 $18{:}00$ 的
集合里再加 $12{:}00$"的两条边，而在这两条边上\textbf{计划费几乎不动}（分别只降
$934$ 与 $808$ 元，其余 $10$ 条边降 $60455$--$354829$ 元）——也就是说两个方案
的信息集\textbf{已经实质等价}；此时增量的大小与符号只能由两次\textbf{独立标定}
的残差决定，而本文恰恰让每种组合各自标定 $\alpha,\rho,\lambda$。把它们读成
"信息已无增量"才是正确的，读成"预报有害"则是把标定噪声当成了因果。

其二，增量对\textbf{费用构成}的作用有清晰的规律。$12$ 条覆盖边上，
\textbf{计划费 $12/12$ 严格下降、调整费 $12/12$ 严格上升}——信息越多，策略越敢
少买常规电量（4-3 全年计划量由 $21931684$~kWh 降至 $21344253$~kWh），把更多
电量留到执行阶段按需追加——但\textbf{紧急费并非同向}，$12$ 条中有 $2$ 条上升。
紧急费是个\textbf{事后}量，取决于缺口落在哪些时段以及当时的价格，而不只是缺口
的大小：八个组合的紧急购电均价在 $5.10$--$5.88$~元/kWh 之间摆动，同样多的缺口
落在不同时段，代价并不相同。这正是本文反复强调的那条原则在评价层面的体现：
判断一版预报值多少钱，只能\textbf{按实际价格结算总费用}，不能只看它减少了多少
购电量或多少次追加。


\subsection{问题四：三种价格信息条件下策略费用的分解图（原正文 \ref{sec:p4-strategy} 节）}

正文表~\ref{tab:p4-strategy} 已给出三种策略在两个分支上的精确费用，这里附上
其费用分解图，便于直观看清"波动电价预测策略把成本从执行阶段挪到计划阶段"
这一结构差异。
\begin{figure}[htbp]
  \centering
  \includegraphics[width=.78\textwidth]{p4_fig3_strategy.png}
  \caption{三种价格信息条件下的策略费用分解（同一实际价格路径结算）}
  \label{fig:p4-strategy}
\end{figure}


\subsection{问题四：$12{:}00$ 与 $18{:}00$ 版本精度倒挂的成因（原正文 \ref{sec:p4-forecast} 节）}
\label{app:p4-inversion}

正文表~\ref{tab:p4-skill} 的"共同时段"一列里，四个发布版本按精度排序为
$6{:}00$（$0.0389$）$<18{:}00$（$0.0403$）$<12{:}00$（$0.0409$）$<0{:}00$
（$0.0428$）。三个更新版本都优于不修正的 $0{:}00$ 版，这符合"日内修正有效"
的预期；但 $12{:}00$ 版反而比 $18{:}00$ 版差 $0.0006$~元/kWh，与"发布越晚、
可用信息越多、判断越准"的直觉相反。这一处不倒饰，下面把它查清。

最容易想到的解释是"$12{:}00$ 的窗口更长（$72$ 段），把上午的偏差带进了
下午"。这个解释是错的，用一个对照就能否掉：把三个版本的 $\beta$ 都换成
同一个值，倒挂立刻消失。
\begin{center}
\footnotesize
\begin{tabular}{lcccc}
  \toprule
  $\beta$ 取法 & $0{:}00$ 版 & $6{:}00$ 版 & $12{:}00$ 版 & $18{:}00$ 版 \\
  \midrule
  冻结值 $(1.0,\,0.75,\,0.5)$ & $0.0428$ & $0.0389$ & $0.0409$ & $0.0403$ \\
  统一 $\beta=0.5$             & $0.0428$ & $0.0389$ & $0.0401$ & $0.0403$ \\
  统一 $\beta=0.75$            & $0.0428$ & $0.0385$ & $0.0409$ & $0.0419$ \\
  \bottomrule
\end{tabular}
\end{center}
统一取 $\beta=0.5$ 时得到 $0.0401<0.0403$，统一取 $\beta=0.75$ 时得到
$0.0409<0.0419$——两种情形下都恢复了"越晚越准"。可见倒挂\textbf{不是窗口
长度造成的}（窗口长度在这三种取法下毫无变化），而是\textbf{逐时刻给
$12{:}00$ 选出的 $\beta=0.75$ 大于给 $18{:}00$ 选出的 $\beta=0.5$}：在末尾
这 $36$ 段上，更强的修正反而吃亏。

这条对照给出一个有用的一般结论：\textbf{逐时刻分别选 $\beta$ 虽然更贴合
各版本的信息量差异，却同时把 $\beta$ 自身的估计噪声带了进来}。选型期的
$\beta$ 扫描对窗口敏感（只看第 $22$--$30$ 天的 $9$ 天，最优值退到
$(0.5,\,0.5,\,0.0)$），正说明这个估计量本身方差偏大。本文因此坚持两个
做法：一是 $\beta$ 只在评价区间之外选一次、全程冻结；二是不为了让
"越晚越准"这一叙事好看而去改 $\beta$——叙事让位于冻结参数的实际结果。

\subsection{问题二、问题三、问题四：六项与七项强制检查的逐条明细}
\label{app:checks}

\subsubsection{六项强制检查}

按前述要求，对 $365$ 天的完整运行逐日执行六项检查，
表~\ref{tab:p2-checks} 给出结果，\textbf{全部通过}：

\begin{enumerate}
  \item \textbf{信息时点}：程序逐日记录预测实际引用的历史日下标，其上界严格
        小于当天；误差分位数只取 $i<n$ 的已结束日期；$g_{n,t}$ 在 $0{:}00$
        冻结，执行层结构中不存在修改 $g$ 的入口。
  \item \textbf{物理可行性}：$365$ 天逐段供需平衡残差的\textbf{最大值}为
        $2.3\times10^{-13}$~kWh（判定阈值为 $1.3\times10^{-3}$~kWh，
        实测值比阈值小 $10$ 个数量级）；储电量全程落在
        $[1200,10800]$~kWh、越界量为 $0$；$c,d\le M$（越限量 $0$）且
        无同时充放电时段（$0$ 段）。
  \item \textbf{跨日连续性}：逐日校验
        $E_{n,0}^{\mathrm{actual}}=E_{n-1,144}^{\mathrm{actual}}$，$365$ 天
        全部成立。
  \item \textbf{状态累计}：式 \eqref{eq:p2-accum} 逐日成立，$\Delta E$ 与
        $\eta_c\sum_t c_t-\eta_d^{-1}\sum_t d_t$ 的最大绝对偏差为
        $1.3\times10^{-11}$~kWh。
  \item \textbf{费用一致性}：用未四舍五入的计划量与紧急量重新计算费用，与
        记录值的\textbf{最大偏差为 $0$ 元}；弃电不退款；日末储备惩罚未混入
        实际账单。
  \item \textbf{汇总一致性}：日期完整、无跨日错位；紧急购电事件合并前后的
        电量总和严格相等（$390$ 个事件全部核对）。
\end{enumerate}


\subsubsection{问题四的六项强制检查（原正文 8.6.2 节）}

对问题四两个分支的 $334$ 个交付日逐日执行下列六项检查，
表~\ref{tab:p4-check} 给出结果，\textbf{未通过天数全部为零}：

\begin{enumerate}
  \item \textbf{信息边界}：第 $\tau$ 次求解实际读取的价格预测与光伏预报，其
        发布时刻均不晚于 $\tau$；$6{:}00$ 之后的价格预测只使用当天
        \textbf{已经实现}的时段，不引用任何未来价格。这一项在 4-2 与 4-3 上
        都按同一标准检查。
  \item \textbf{计划冻结（仅 4-2）}：$g_t$ 在整个交付日内恒等于 $0{:}00$ 定下的
        $g_t^0$，执行层不存在修改 $g$ 的入口。4-3 不受这一约束，改由下一条
        检查它的调整权限。
  \item \textbf{调整权限（仅 4-3）}：$x_t\ge g_t^0$ 在全部时段成立，主计费口径下
        下调量恒为零；交付块边界处 $x_t$ 的取值只由该块发布时刻的信息决定，
        不跨块引用更晚的信息。
  \item \textbf{物理可行性}：$334$ 天逐段供需平衡残差的\textbf{最大值}为
        $2.3\times10^{-13}$~kWh（两个分支相同）；储电量全程落在
        $[1200,10800]$~kWh，\textbf{上下界余量均为 $0$}（即恰好触界而不越界）；
        不存在同时充放电的时段。
  \item \textbf{状态累计}：由 $c_t,d_t$ 从日初储电量独立递推得到的序列与记录值
        一致；当日日末电量严格作为次日初值，\textbf{跨日衔接最大偏差为
        $0$~kWh}。
  \item \textbf{费用一致性}：计划费、调整费、紧急费三项分别按实际价格重算，
        与记录值的\textbf{最大偏差均为 $0$ 元}；$0{:}00$ 的"预估费用"另按当时
        的预测价复算，同样相符。这正是第 \ref{sec:p4-model} 节所强调的原则：
        \textbf{决策用当时的信息，结算一律用实际价格}。
\end{enumerate}


\subsection{问题一：最优解六项复核的逐条明细（原正文 5.3.3 节）}
\label{app:p1-checks}

\subsubsection{求解结果的校验}

为确保结果可信，对最优解独立复核了以下各项，\textbf{全部通过}：

\begin{enumerate}
  \item \textbf{时段完整性}：$144$ 个时段完整覆盖当天 $0{:}00$--$24{:}00$，
        首尾与四小时分组均无错位；附件 1 时间标签经校验为
        $\{10,20,\dots,1440\}$ 分钟且严格递增。
  \item \textbf{逐段供需平衡}：$\max_t\bigl|g_t+v_t-s_t+d_t-\ell_t-c_t\bigr|
        < 10^{-6}$~kWh。
  \item \textbf{状态递推}：由 $c^*,d^*$ 从 $E_0$ 出发独立累计递推得到的
        储电量序列与求解器返回的 $E^*$ 最大偏差小于 $10^{-6}$~kWh。
  \item \textbf{物理边界}：$E^*$ 全程落在 $[1200,10800]$~kWh 内；
        $c_t^*,\,d_t^* \le M = 2500/3$~kWh；$0 \le s_t^* \le v_t$；
        且不存在同时充放电的时段。
  \item \textbf{日循环}：$E_{144}^* = E_0^* = 6000$~kWh，且按式
        \eqref{eq:cycle-check} 校验 $\sum_t d_t^* = 0.81\sum_t c_t^*$
        （实测 $16799.94 / 20740.67 = 0.810000$）。
  \item \textbf{费用与汇总}：用未四舍五入的数据重新计算
        $\sum_t p_t g_t^*$，与求解器返回值一致；各分段汇总与全天明细一致。
\end{enumerate}

数值复核仅设置极小的浮点容差，\textbf{未通过删除负数、裁剪越界值或修改末段
结果的方式掩盖任何约束违背}。

\subsection{问题二：四个指定日期的逐日结构（原正文 6.6.3 节）}
\label{app:p2-dates}

\subsubsection{指定日期的结果}

表~\ref{tab:p2-buy}、表~\ref{tab:p2-storage}、表~\ref{tab:p2-storage-dec} 与
表~\ref{tab:p2-emg} 分别给出 $3$ 月 $20$ 日、$6$ 月 $21$ 日、$9$ 月 $23$ 日与
$12$ 月 $21$ 日四个指定日期的计划购电量、储能充放电量与紧急购电情况（其中
$12$ 月 $21$ 日的储能充放电量因列宽原因单列一表）。四个日期的运行特征差异
明显，恰好覆盖了不同季节与不同星期的组合：

\begin{itemize}
  \item \textbf{$3$ 月 $20$ 日（周四）}与\textbf{$9$ 月 $23$ 日（周二）}为
        工作日，负载处于高位（$117648.38$~kWh 与 $121182.28$~kWh），光伏
        出力中等（$55439.42$~kWh 与 $59065.67$~kWh），计划购电量分别为
        $67400.98$~kWh 与 $66102.01$~kWh；
  \item \textbf{$6$ 月 $21$ 日（周六）}属于低负载的周末族，全天负载仅
        $81408.06$~kWh，而光伏出力却是四者中最高（$62073.32$~kWh），
        净负荷最低，故计划购电量仅 $40911.64$~kWh，为四个日期中最低；也是
        四者中弃电量最大的一天（$19087.44$~kWh）；
  \item \textbf{$12$ 月 $21$ 日（周日）}适逢冬至，光伏出力处于全年低位
        （$35338.66$~kWh，在正式区间内为第 $8$ 低），负载则属高负载族
        （$124376.44$~kWh，该族约 $124000$~kWh/日），净负荷在正式区间内
        列第 $3$，全天计划购电量高达 $99337.46$~kWh、购电费
        $64146.24$~元，为四个日期中最高。
\end{itemize}

四个日期的这一组差异，恰好覆盖了负载的星期结构与光伏的季节变化两个维度，
使得指定日期的结果具有代表性。

在紧急购电方面，四个指定日期中有两个\textbf{全天无紧急购电}（$6$ 月 $21$ 日
与 $12$ 月 $21$ 日），$3$ 月 $20$ 日仅 $1$ 个事件（$20{:}30$--$20{:}40$，
$14.56$~kWh、$101.56$~元），而 $9$ 月 $23$ 日发生 $3$ 个紧急购电事件
（$18{:}50$--$20{:}40$、$20{:}50$--$22{:}00$、$22{:}20$--$22{:}30$），合计
$1253.46$~kWh、$7777.51$~元。这些事件都出现在
晚间 $18{:}50$ 以后：该日白天光伏出力偏低，储能到傍晚已接近下限，晚间高价时段
出现预测未能覆盖的额外缺口，只能按 $5$ 倍电价补购。事件按"同日连续发生紧急
购电的时段合并为一个事件"的规则生成，合并后的电量与逐段之和严格相等。

\subsection{问题三：调整次数、储能轨迹与逐日收益（原正文 7.7 节（3）--（5））}
\label{app:p3-analysis}

\textbf{（3）调整并非越多越好。}三个发布时刻全年共正式提交
$334\times3=1002$ 次，但其中相当一部分提交的结果是"保持不变"：全部
$334\times144=48{,}096$ 个时段中，最终生效量被真正改动的只有 $16{,}788$ 个，
占 $34.91\%$。图~\ref{fig:p3-purchase} 给出四个指定日期的初始计划购电量与最终
生效购电量，两条曲线的差异只在少数时段出现。每个交付块的调整都是\textbf{双向}的
（全年上调 $656{,}263.05$~kWh、调减 $775{,}660.56$~kWh），这与问题二的执行规则
形成了对照——问题二只能被动接受计划、按 $5$ 倍电价补缺。这也说明更新预报后重新
求解\textbf{并不必然改变购电量}：既可以调整购电，也可以只重排储能；只有当预计
减少的紧急购电收益超过调整费时，购电量才真正被改动。

\textbf{（4）储能轨迹被滚动修正。}图~\ref{fig:p3-storage} 给出四个指定日期的
参考轨迹与实际储电量。可以看到每个发布时刻的参考轨迹都会重新拼接（竖虚线），
而实际储电量始终连续。评价区间内\textbf{日末}储电量落在 $5{,}216.7$--$7{,}966.9$
kWh 之间，均值 $6{,}284.9$~kWh，没有出现日末被放空或充满的极端状态。这里必须
区分\textbf{日末位置}与\textbf{逐时轨迹}：逐时段看，储能确实会在光伏富余时段
触及 $10800$~kWh 的上界（$334$ 天中有 $322$ 天）、在连续缺口时被压至
$1200$~kWh 的下界（$146$ 天），这正是设备容量对消纳能力的真实约束；而储备线
$R_t=\mathrm{clip}(1200+\rho(\bar E_t-1200),1200,10800)$ 的作用，是防止执行层
为省钱而在日末透支储能。四个指定日期的具体结果见表~\ref{tab:p3-buy}、
表~\ref{tab:p3-storage}、表~\ref{tab:p3-storage-dec} 与表~\ref{tab:p3-emg}。

\begin{figure}[htbp]
  \centering
  \begin{subfigure}{\textwidth}
    \centering
    \includegraphics[width=.66\textwidth]{p3_fig2_purchase.png}
    \caption{逐小时汇总的初始计划购电量 $g_t^0$ 与最终生效购电量 $x_t$}
    \label{fig:p3-purchase}
  \end{subfigure}\\[3pt]
  \begin{subfigure}{\textwidth}
    \centering
    \includegraphics[width=.66\textwidth]{p3_fig3_storage.png}
    \caption{滚动参考储电量 $\bar E_t$ 与实际储电量 $E_t$}
    \label{fig:p3-storage}
  \end{subfigure}
  \caption{指定日期的购电量调整与储能轨迹}
\end{figure}

\textbf{（5）收益在月份间并不均匀。}图~\ref{fig:p3-daily} 给出主策略逐日的
调整费、紧急费与日末储电量。按月汇总，\textbf{调整费在 $12$ 月最高}
（$164.18$~万元），约为 $4$ 月（$90.60$~万元）的 $1.81$ 倍；\textbf{紧急费则在
$5$ 月最高}（$5.61$~万元）、$11$ 月最低（$0.89$~万元），跨度接近 $6$ 倍。调整费
的峰值出现在年末，与冬季光伏出力低、净负荷高、计划量对误差更敏感有关；紧急费的
峰值出现在 $5$ 月，与空调负荷开始爬升、日内波动加大有关。两项支出的月度分布并
不同步，这也说明分月评价预报收益是必要的，不能只看四个指定日期。

\begin{figure}[htbp]
  \centering
  \includegraphics[width=.66\textwidth]{p3_fig5_daily.png}
  \caption{主策略逐日的调整费、紧急费与日末储电量}
  \label{fig:p3-daily}
\end{figure}
